This dissertation adds an uncertainty quantification layer to a reconstructed continuous field on a sphere from sparse, scattered observations. This is relevant to areas such as climate science, geophysics, and computer graphics, where uncertainty is important for real-world applications.

From earlier research, the physics-informed spherical harmonic model developed by Thomas (2025) uses a regularised spherical harmonic expansion with a Sobolev-type seminorm penalty to produce accurate point reconstructions of global temperature from city-level data. However, it provides no information about the reliability of these predictions. 

Therefore, adaptive geodesic conformal prediction was applied as an additional layer to the existing spherical harmonic pipeline. Since temperature is scalar rather than manifold-valued, spherical geometry enters through the input locations rather than the output. By using great-circle distances on the sphere, a cross-validated nearest-neighbour estimator was used to estimate the local reconstruction difficulty, $\hat{\sigma}(x)$. The conformal scores are then normalised by $\hat{\sigma}(x)$ before calibration so that the resulting prediction interval width varies according to local data density and local reconstruction difficulty.

This project used Berkeley Earth city-temperature observations, with $n=3,510$ observations from July 2003. The standard conformal prediction method achieved an overall coverage close to the nominal 90\%; however, significant differences were observed across difficulty levels. The lowest-difficulty tertile achieved 100\% coverage, whereas the highest-difficulty tertile achieved only 76.7\%. In contrast, the adaptive method produced more balanced coverage across three tertiles, between 91\% and 95\%. The adaptive method produced a narrower mean prediction interval of 4.59°C, compared with 4.65°C for the standard method. Overall, the adaptive method provides more reliable local uncertainty estimates while maintaining a similar mean interval width. 

A robustness experiment across spherical harmonic degrees from L=8 to L=30 showed that the adaptive method continued to perform well across different model complexities. The highest-uncertainty locations were mainly concentrated in the high-altitude Andes, where local observation coverage was relatively sparse. This may indicate that the estimated uncertainty captures aspects of spatial data sparsity and provides useful insight into where reconstruction predictions may be less reliable.

\noindent \textbf{Keywords:} spherical harmonics; conformal prediction; uncertainty quantification; climate field reconstruction; geodesic distance; adaptive calibration; spatial statistics; Sobolev regularisation
